% Chapter 16 -- Claim: Selling Options Earns a Reliable Premium
% Requires: amsmath, amssymb, booktabs, array, tabularx, graphicx, tcolorbox (breakable, skins), hyperref or url.
% Defines the callout box only if the preamble has not.
\providecolor{fenavy}{HTML}{1B2A6B}
\providecolor{feshade}{HTML}{F3F5FA}
\makeatletter
\@ifundefined{callout}{%
  \newtcolorbox{callout}{breakable, enhanced,
    colback=white, colframe=gray!70, boxrule=0.5pt, arc=1mm,
    borderline west={2.5pt}{0pt}{gray!70},
    left=8pt, right=8pt, top=4pt, bottom=4pt}%
}{}
\makeatother

\chapter{Claim: Selling Options Earns a Reliable Premium}
\label{ch:16}

Once the Black--Scholes model gave every option an implied volatility, a regularity became visible: the volatility implied by index option prices was usually higher than the volatility the market went on to deliver. An investor who sells stock-index options is in the same business as an insurance company. A put option\index{option!put} is insurance against a market fall; its seller collects a premium up front and pays out when the market drops. If the insurance is priced above expected claims, the seller earns a steady profit, with the occasional very bad year. Coval and Shumway (2001) showed that index option returns are too low to be explained by their market risk, and Carr and Wu (2009) measured the \textbf{variance risk premium}\index{variance risk premium} directly and found it large and persistent.

The question matters because a large industry rests on it. Whaley (2002) studied the covered-call\index{option selling!covered call} strategy and helped design the Cboe BuyWrite Index, and option-income funds, covered-call ETFs and volatility-selling strategies now hold large sums. The issue is what kind of premium they collect. A reliable premium for bearing a risk that others want to shed is one thing. A strategy that wins small most months and loses large in a few, while carrying most of the stock market's risk, is another. Israelov and Nielsen (2015) showed that much of the return of covered calls comes from their equity exposure rather than from the volatility premium.

The implications are practical for anyone offered option income. The two kinds of premium can look alike in a backtest\index{backtest}, so the evidence has to separate them, and the crashes of 2008, February 2018 and March 2020 show what the tail looks like. The claim on trial\index{claims-on-trial method} is: \textbf{selling index puts or writing covered calls earns a return above the risk taken, because implied volatility exceeds realized volatility.}

The lab in Section~\ref{sec:16-lab} replicates the variance-risk-premium measurement of Carr and Wu (2009) with the free VIX and daily U.S. market returns, and rebuilds put-writing and covered-call strategies that can be compared with Cboe's published PUT and BXM index histories, in order to check whether selling options earns a reliable premium and how much of it is simply stock-market risk.

By the end of the chapter you will be able to:

\begin{itemize}
  \item describe call and put payoffs and use put--call parity\index{put--call parity} to show that a covered call and a cash-secured put\index{option selling!cash-secured put} are the same position;
  \item price an option by Black--Scholes\index{Black--Scholes--Merton model}, compute its Greeks\index{Greeks}, and back out implied volatility\index{volatility!implied};
  \item define the VIX\index{VIX} and the variance\index{variance} risk premium\index{risk premium}, and estimate a seller's expected gain from the gap between implied and realized volatility\index{volatility!realized};
  \item state the claim that selling options earns a reliable premium, and separate the variance risk premium from the returns of option-income strategies;
  \item measure option-strategy performance with beta\index{beta}-matched benchmarks, up and down betas\index{beta!up and down} and the manipulation-proof performance measure\index{manipulation-proof performance measure (MPPM)};
  \item judge what an option-income fund is likely to deliver, and what can go wrong.
\end{itemize}

\section{Foundations: the VIX and the option seller's payoff}
\label{sec:16-foundations}

Chapter~\ref{ch:14} introduced option payoffs, put--call parity, the Black--Scholes--Merton model, the Greeks and implied volatility (Sections~\ref{sec:14-bsm} to~\ref{sec:14-iv}). This chapter adds the VIX and the economics of selling options.

\subsection{The VIX}

The VIX, introduced in Section~\ref{sec:10-iv}, is computed from a strip of out-of-the-money calls and puts on the S\&P 500. Its square approximates the fair strike of a 30-day \textbf{variance swap}\index{variance swap|textbf}, a contract that pays the difference between realized variance and a fixed strike. The variance risk premium of Section~\ref{sec:10-iv} is therefore the expected payoff to the seller of such a swap: the expected gap between implied (risk-neutral\index{risk-neutral pricing}) variance and the realized variance that follows.

\textbf{Worked example.} The index is at 100, the risk-free rate\index{risk-free rate} is 4\%, and a one-month at-the-money put trades at an implied volatility of 20\%.

\begin{itemize}
  \item Black--Scholes price: $P = \$2.14$ per \$100 of index. The call is worth \$2.47; the difference, \$0.33, equals $100 - 100\,e^{-0.04/12}$, as parity requires.
  \item Suppose volatility over the month turns out to be 16\%. At 16\%, the put would have been worth \$1.68. The seller's expected gain is roughly $\$2.14 - \$1.68 = \$0.46$, or vega (\$0.115 per vol point) times the 4-point gap.
  \item That is an \emph{expected} gain. If the index falls 10\% during the month, the put pays \$10 and the seller loses $\$10 - \$2.14 = \$7.86$, about seventeen months of expected gains at \$0.46 a month (Figure~\ref{fig:16-put}).
\end{itemize}

\begin{figure}[htbp]
\centering
\begin{tikzpicture}[declare function={profit(\s) = 2.14 - max(100 - \s, 0);}]
\begin{axis}[
  width=0.64\linewidth, height=0.34\linewidth,
  axis lines=left, axis line style={-stealth},
  xmin=85, xmax=110, ymin=-13, ymax=5,
  xtick={86,88,...,110}, ytick={-10,-5,0,5},
  tick label style={font=\small},
  xlabel={Index at expiry}, ylabel={Put seller's profit (\$)},
  clip=false,
]
  \addplot[gray, thin, domain=85:110, samples=2] {0};
  \addplot[fenavy, very thick] coordinates {(85,-12.86) (100,2.14) (110,2.14)};
  \fill[fenavy] (axis cs:97.86,0)  circle (2.6pt);
  \fill[fenavy] (axis cs:90,-7.86) circle (2.6pt);
  \node[anchor=south east] at (axis cs:97.4,0.4) {break-even 97.86};
  \node[anchor=south east] at (axis cs:110,2.4) {keep the \$2.14 premium};
  \node[anchor=north west, align=left] at (axis cs:90.6,-7.3)
    {index $-10\%$: lose \$7.86,\\about 17 months of expected gains};
\end{axis}
\end{tikzpicture}
\caption{The put seller's profit at expiry in the worked example. Above the strike of 100 the seller keeps the \$2.14 premium; below the break-even of 97.86 the loss grows one for one with the fall. A 10\% fall costs \$7.86, about seventeen months of expected gains.}
\label{fig:16-put}
\end{figure}

\section{How it works}
\label{sec:16-how}

\textbf{Why insurance might be overpriced.} Many investors hold stock portfolios they cannot easily sell and want protection against crashes. Pension funds, insurers and structured-product issuers buy index puts; market makers\index{market maker} who sell them must hold capital against the risk and cannot fully hedge it. Bollen and Whaley (2004) and G\^arleanu, Pedersen and Poteshman (2009) showed that this net buying pressure raises option prices, especially for out-of-the-money puts. In equilibrium, sellers are paid for bearing a risk that bites in the worst states of the world, when investors' wealth is low and losses hurt most. A premium for that risk is not a free lunch. It is compensation, and it should be large precisely because the losses come at bad times.

\textbf{Ways to sell volatility.} The strategies differ in how much equity risk they carry alongside the volatility exposure.

\begin{center}
\small
\begin{tabularx}{0.9\linewidth}{@{}>{\raggedright\arraybackslash}p{0.19\linewidth}>{\raggedright\arraybackslash}X>{\raggedright\arraybackslash}X>{\raggedright\arraybackslash}X@{}}
\toprule
Strategy & Position & Equity exposure & Main risk \\
\midrule
Covered call (buy-write) & Own the index, sell a call & High (beta about 0.5--0.8) & Market falls; gains capped in rallies \\
Cash-secured put (put-write) & Hold bills, sell a put & High (same as covered call at the same strike) & Market falls \\
Short straddle & Sell a call and a put & Low on average, changes with the market & Large moves either way \\
Delta-hedged option & Sell an option, hedge its delta daily & Near zero & Realized volatility above implied \\
Short variance swap & Receive fixed variance, pay realized & Near zero, but crashes raise realized variance & Volatility spikes \\
\bottomrule
\end{tabularx}
\end{center}

The first two are what most retail and institutional ``option income'' products do. The last two isolate the variance risk premium itself. Comparing them is the heart of this chapter: if the covered call's return is mostly the equity premium\index{equity premium} earned with a beta of 0.6, it is not evidence that selling options pays.

\textbf{The payoff shape.} A covered call keeps the downside of the stock market but caps the upside at the strike plus the premium (Figure~\ref{fig:16-cc}). In calm or gently rising months it beats the market; in strong rallies it lags; in crashes it loses almost as much as the market, cushioned only by the premium. The return distribution is therefore \textbf{negatively skewed}\index{skewness|textbf}. Most months bring small gains, and a few bring large losses.

\begin{figure}[htbp]
\centering
\begin{tikzpicture}
\begin{axis}[
  width=0.61\linewidth, height=0.38\linewidth,
  axis lines=left, axis line style={-stealth},
  xmin=80, xmax=120, ymin=-20, ymax=21,
  xtick={80,85,...,120}, ytick={-20,-10,0,10,20},
  xlabel={Index at expiry}, ylabel={Profit per \$100 of index (\$)},
  clip=false,
]
  \addplot[gray, thin, domain=80:120, samples=2] {0};
  % Index alone
  \addplot[fegray, thick, dashed, domain=80:120, samples=2] {x - 100};
  % Covered call: index plus the $2.47 premium, capped at the strike
  \addplot[fenavy, very thick] coordinates {(80,-17.53) (100,2.47) (120,2.47)};
  \node[fegray, anchor=east] at (axis cs:113,16.5) {index alone};
  \node[fenavy, anchor=south east, align=right] at (axis cs:120,3.0)
    {covered call:\\gain capped at \$2.47};
  \node[fenavy, anchor=west, align=left] at (axis cs:90,-15)
    {losses nearly as large\\as the index's};
\end{axis}
\end{tikzpicture}
\caption{A one-month at-the-money covered call against the index alone, using the worked example's prices. The call premium of \$2.47 caps the gain and cushions the loss by the same amount, so the covered call keeps almost all of the index's downside.}
\label{fig:16-cc}
\end{figure}

That shape creates a measurement problem. The Sharpe ratio uses only the mean and standard deviation, so it cannot see skewness. Selling out-of-the-money options can raise a portfolio's Sharpe ratio while making its tail much worse. Section~\ref{sec:16-measure} sets out measures that are harder to game.

\section{The claim}
\label{sec:16-claim}

\begin{callout}
\textbf{Claim.} Implied volatility exceeds the volatility that follows, on average, so option sellers earn a reliable premium. Strategies that sell index options, such as covered calls and put-writing, deliver better risk-adjusted returns than holding the index.
\end{callout}

The claim has two parts that can succeed or fail separately.

\begin{itemize}
  \item \textbf{The variance risk premium.} Implied variance is, on average, above the realized variance that follows. This is a statement about prices and can be tested directly by comparing the VIX with subsequent realized volatility.
  \item \textbf{The strategy premium.} Covered-call and put-writing strategies earn more than their risk justifies. This requires the variance premium to survive trading costs, early exercise and roll timing, and it requires an honest risk adjustment.
\end{itemize}

To test it we need to be precise on three points.

\begin{itemize}
  \item \textbf{What counts as risk.} A strategy with a beta of 0.5 should be compared with half the market plus half in bills, not with the whole market. Its negative skewness should be priced, not ignored.
  \item \textbf{What counts as reliable.} A premium that is positive in most months but wiped out in a few is reliable in one sense and not in another. We report both the share of months positive and the size of the worst losses.
  \item \textbf{What would count against it.} The variance risk premium fails if implied volatility is not significantly above realized volatility on average. The strategy claim fails if covered calls and put-writing earn no alpha after adjusting for beta, or if their advantage disappears once skewness is priced.
\end{itemize}

\section{The original evidence}
\label{sec:16-evidence}

\textbf{Option returns.} Coval and Shumway (2001) studied S\&P 500 index options and found that buying zero-beta, at-the-money \textbf{straddles} (a call and a put at the same strike) lost about 3\% a week on average. A straddle has little exposure to the direction of the market, so its losses point to a price for volatility itself. Bakshi and Kapadia (2003) found that delta-hedged option positions lost money on average, and attributed the losses to a negative market volatility risk premium: investors pay to be long volatility because volatility rises when markets fall.

\textbf{The variance risk premium.} Carr and Wu (2009) showed how to compute the variance swap rate from a portfolio of options, the method behind today's VIX, and used it to measure the premium directly for five stock indexes and 35 individual stocks. For the S\&P 500, implied variance was well above realized variance on average. Bollerslev, Tauchen and Zhou (2009) found that the size of the premium also predicts stock-market returns at a quarterly horizon, better than the price--earnings ratio or the default spread, consistent with its being a price of economic uncertainty.

\textbf{Covered calls and put-writing.} Cboe commissioned Robert Whaley to design the BuyWrite Index (BXM), which tracks an S\&P 500 portfolio that sells a one-month, near-the-money call each month. Whaley (2002) found that it had earned returns close to the S\&P 500's with substantially lower volatility. A later study for Cboe (Callan Associates 2006) reported that from June 1988 to August 2006 the BXM returned 11.77\% a year against 11.67\% for the S\&P 500, with a standard deviation of 9.29\% against 13.89\%. Cboe later added a PutWrite Index (PUT), which sells one-month at-the-money puts against Treasury-bill collateral. These results drew large inflows into option-income products.

\textbf{The skeptics.} Broadie, Chernov and Johannes (2009) argued that large average losses on index puts are not by themselves evidence of mispricing. Under standard models with jumps and reasonable risk premiums, put returns are expected to be very negative and very noisy, so short samples say little. Israelov and Nielsen (2015) decomposed covered-call returns into three exposures: the equity market, short volatility and an unintended equity reversal bet. Equity exposure drove most of the risk and return. The short-volatility component had a Sharpe ratio near 1.0 but contributed less than 10\% of the risk, and the reversal component added about a quarter of the risk with little return. A covered call, in other words, is mostly a stock-market position.

\section{Measuring option-strategy returns}
\label{sec:16-measure}

Option-selling strategies are the textbook case where the usual performance measures mislead. Four corrections matter.

\textbf{Adjust for beta.} A covered call with a beta of 0.6 should earn about 60\% of the equity premium without any volatility premium at all. The relevant comparison is a \textbf{beta-matched} portfolio, 60\% in the market and 40\% in bills, or equivalently the CAPM alpha (Chapter~\ref{ch:18}). A covered call that matches the market's return with two-thirds of its volatility may have earned no alpha.

\textbf{Look at both betas.} Selling options makes exposure asymmetric. The \textbf{down beta}, estimated in months when the market falls, is close to 1: the strategy takes most of the loss. The \textbf{up beta}, estimated in rising months, is much lower, because the call caps the gains. An investor who cares about losses more than gains should weigh the down beta.

\textbf{Price the skewness.} The Sharpe ratio ignores skewness, so a strategy can raise it by selling crash insurance. Goetzmann, Ingersoll, Spiegel and Welch (2007) showed that a manager can raise a portfolio's Sharpe ratio by selling out-of-the-money options, without any skill. They proposed a \textbf{manipulation-proof performance measure (MPPM)}:
\begin{equation}
  \Theta = \frac{1}{(1 - \rho)\,\Delta t}\,\ln\!\left\{\frac{1}{T}\sum_{t=1}^{T}\left[\frac{1 + r_t}{1 + r_{f,t}}\right]^{1-\rho}\right\}
  \label{eq:16-mppm}
\end{equation}
Here $T$ is the number of periods, $\Delta t$ the length of each period in years, and $r_t$ and $r_{f,t}$ the strategy's and the risk-free return. $\Theta$ is the certainty-equivalent excess return of an investor with constant relative risk aversion $\rho$. With $\rho = 3$, a typical value, a large negative month costs far more than a large positive month gains, so selling crash insurance cannot improve the score unless it really adds value.

\textbf{Respect the sample.} A strategy that loses heavily once a decade needs several decades of data to estimate its mean (Section~\ref{sec:01-foundations}). Broadie, Chernov and Johannes (2009) showed how noisy put returns are even over twenty years. A backtest that starts after 2009 and ends before 2020, for example, would miss both of the big losses of the period since 1990.

\section{Lab~\thechapter: Replicating Carr and Wu (2009) and the Cboe put-write and buy-write indices}
\label{sec:16-lab}

\begin{callout}
\textbf{Lab~\ref{ch:16} at a glance.} \emph{Papers replicated:} the variance risk premium of Carr and Wu (2009); the put-write and buy-write strategies studied by Whaley (2002). \emph{Claim on trial:} selling index options earns a reliable premium. \emph{Data:} daily market returns from the Kenneth French library and the VIX from FRED, which starts in 1990; both are free. Cboe's PUT and BXM index histories can be added for comparison. \emph{Code:} runs as a script (\texttt{python lab14.py}) or as a notebook with a prediction prompt before each step.
\end{callout}

\textbf{Part A, the variance risk premium.} At each month end, the VIX is compared with the market's realized volatility over the next 21 trading days. The lab reports the average gap in volatility points with a Newey--West $t$-statistic, the share of months in which implied exceeded realized, the median gap and the worst months (Lab Figures~\ref{ch:16}.1 and~\ref{ch:16}.2).

\textbf{Part B, strategies.} The lab rebuilds two strategies from the data. Each month it sells a one-month option priced by Black--Scholes at the VIX and holds it to the next month end:

\begin{itemize}
  \item an \textbf{at-the-money put-write}, collateralized by bills like Cboe's PUT. By put--call parity this is also an at-the-money covered call, and the lab checks that the two returns are identical;
  \item a \textbf{covered call 2\% out of the money}, which keeps more of the market's upside.
\end{itemize}

Each is compared with the market on excess return, volatility, Sharpe ratio, skewness, worst month, drawdown, beta and CAPM alpha, over its own history and over a common period (Lab Figure~\ref{ch:16}.3).

\textbf{Part C, robustness checks.} Returns in crisis periods, the five worst market months, up and down betas, and the MPPM (Lab Figure~\ref{ch:16}.4).

\textbf{Part D, equity versus volatility.} Each strategy's excess return is split into the part explained by beta and the alpha that remains. The lab compares it with the market scaled to the same beta, and with a pure short variance swap struck at VIX$^2$, sized to 10\% annual volatility (Lab Figure~\ref{ch:16}.5).

The rebuilt strategies are deliberately simple, and that matters for interpretation. They price options at the VIX, which includes the premium on out-of-the-money puts and so usually exceeds at-the-money implied volatility. That overstates the premium a seller actually collects. They also ignore bid--ask spreads\index{bid--ask spread} and roll at month ends rather than on Cboe's third-Friday schedule. The \texttt{-{}-iv-shift} option prices the options a few points below the VIX. Comparing the rebuilt series with Cboe's indices, which use traded prices, shows how much these simplifications matter.

\textbf{What the simulated run shows.} The simulation builds in a VIX 12\% above expected volatility. The lab finds the premium (implied above realized in 83\% of months), a put-write with a beta of about 0.5, strongly negative skewness and a down beta near 0.9 against an up beta near 0.15. The rebuilt put-write and covered call show positive alphas\index{alpha} when options are priced at the inflated VIX. The Cboe-like series, priced 2 points lower, show none. These results confirm that the code works. They say nothing about real markets. Run the lab on the real data to reach your own verdict.

\subsection{Robustness checks}
\label{sec:16-stress}

The full-sample results can hide weaknesses. The checks below test the claim under the conditions in which it is hardest to meet.

\textbf{Crashes.} The months that decide the outcome are few: October 1987, August 1998, September 2008 to February 2009, February and March 2020. In each, the market fell sharply, realized volatility jumped far above the implied volatility priced a month earlier, and option sellers lost most of what the market lost. A put-write cushions a crash only by one month's premium, typically 1--3\% of the strike. Lab Figure~\ref{ch:16}.2 shows the pattern for the variance premium: many small positive months and a few large negative ones.

\textbf{Skewness.}\index{skewness|textbf} On the Sharpe ratio\index{Sharpe ratio}, put-writing and covered calls often look better than the market, because they cut volatility more than they cut returns. On the MPPM with risk aversion\index{risk aversion} 3, expect the advantage to shrink, because the measure prices the crash months. A strategy whose apparent edge vanishes once skewness is priced is being paid for crash risk\index{crash risk}, not earning alpha.

\textbf{Equity exposure.} Covered calls and put-writes have betas of roughly 0.5 to 0.7. Over a period when the equity premium was large, a beta of 0.6 alone earns most of their return. The lab's decomposition separates the two: if the alpha is small relative to the beta component, the strategy is mainly a stock investment with a different payoff shape.

\textbf{Timing of the sample.} After a crisis, implied volatility can stay high while realized volatility subsides, which makes the premium look rich; before a crisis, it can look rich until the crash arrives. A backtest that starts just after a crisis, or ends just before one, will flatter the strategy. Exercise 4 splits the sample at 2008.

\textbf{Implementation.} Real options have bid--ask spreads, and covered-call funds pay management fees. Israelov and Nielsen (2015) found that the strategy's equity-reversal exposure, created by mechanically resetting the strike each month, adds risk without return. These frictions are small per trade but large relative to a premium of a few percent a year.

\subsection*{Your verdict}

Reach your own verdict on the claim from the numbers in your own run, using the scale of Section~\ref{sec:01-method}. Then compare it with the published evidence summarized in Section~\ref{sec:16-verdict}. If your verdict differs, explain which of your results led you there.

\section{Where the evidence stands}
\label{sec:16-verdict}

The published studies cited in this chapter point to a clear pattern: \textbf{implied volatility has exceeded realized volatility on average, so the variance risk premium is real, but covered-call and put-writing strategies earn most of their return from stock-market exposure and lose most of what the market loses in crashes.} The table summarizes where the evidence supports the claim and where it does not, as of November 2026, and names the source of each finding. The lab lets you verify these findings with free data.

\begin{center}
\small
\begin{tabularx}{\linewidth}{@{}>{\raggedright\arraybackslash}p{0.19\linewidth}>{\raggedright\arraybackslash}p{0.10\linewidth}>{\raggedright\arraybackslash}X>{\raggedright\arraybackslash}p{0.23\linewidth}@{}}
\toprule
Situation & Holds? & Why & Evidence \\
\midrule
Implied above realized variance & \textbf{Yes} & Implied variance has been well above realized variance on average for the S\&P 500 and other indexes. & Carr and Wu (2009) \\
\addlinespace[3pt]
Option buyers pay for volatility & \textbf{Yes} & Zero-beta straddles lost about 3\% a week, and delta-hedged option positions lost money on average. & Coval and Shumway (2001); Bakshi and Kapadia (2003) \\
\addlinespace[3pt]
An economic reason for the premium & \textbf{Yes} & Demand for crash insurance raises option prices, and the size of the premium predicts market returns. & Bollen and Whaley (2004); G\^arleanu, Pedersen and Poteshman (2009); Bollerslev, Tauchen and Zhou (2009) \\
\addlinespace[3pt]
Covered calls against the index & \textbf{Partly} & The BuyWrite Index matched the S\&P 500's return with lower volatility from 1988 to 2006. & Whaley (2002); Callan Associates (2006) \\
\addlinespace[3pt]
Covered calls after adjusting for beta & \textbf{Weakly} & Equity exposure drove most of the risk and return; the short-volatility part was under 10\% of the risk. & Israelov and Nielsen (2015) \\
\addlinespace[3pt]
Risk adjustment that prices skewness & \textbf{Weakly} & Selling options can raise a Sharpe ratio without skill; measures that price the tail show a smaller advantage. & Goetzmann, Ingersoll, Spiegel and Welch (2007) \\
\addlinespace[3pt]
Mispricing, judged from put returns & \textbf{No} & Very negative, noisy put returns are expected under standard models with jumps, so they are not by themselves evidence of mispricing. & Broadie, Chernov and Johannes (2009) \\
\bottomrule
\end{tabularx}
\end{center}

The two parts of the claim come apart. The premium for bearing volatility risk is real; the strategies sold to collect it are mostly stock-market positions with a worse tail.

\textbf{What would change the picture.} The evidence for option-income strategies would strengthen if their alpha against a beta-matched benchmark remained significant after pricing skewness and costs over a full cycle of crises. It would weaken further if, after accounting for beta, the premium were consistently eaten by costs and fees.

\subsection{Open questions}
\label{sec:16-open}

\begin{itemize}
  \item \textbf{Has the boom in short-dated options changed the premium?} Options expiring the same day (zero days to expiry, or 0DTE) made up over 61\% of S\&P 500 index option volume in May 2025, according to Cboe (Cboe Global Markets 2025). Do they change the price of one-month volatility, or the size of the variance premium?
  \item \textbf{How large is the premium net of costs for retail investors?} Covered-call ETFs on single stocks and indexes have grown rapidly. What do they deliver after fees, spreads and taxes\index{taxes}, compared with a beta-matched index position?
  \item \textbf{Is the premium compensation or mispricing\index{mispricing}?} Rational models with crash risk explain much of it (Broadie, Chernov and Johannes 2009). Does demand pressure from hedgers add a mispricing on top?
  \item \textbf{Can the premium be harvested with less crash risk?} Delta-hedged\index{delta hedging} positions, variance swaps with caps, or strategies that cut exposure when volatility rises all aim to keep the premium and avoid the tail. Do they survive out of sample\index{in-sample versus out-of-sample}?
\end{itemize}

\section{What practitioners and investors can use}
\label{sec:16-investors}

The evidence in Section~\ref{sec:16-verdict} suggests several practical lessons.

\textbf{What an option-income fund actually delivers.} Covered-call and put-writing funds advertise high distributions. Three facts help interpret them; Israelov and Nielsen (2014) discuss these and other common misconceptions about covered calls.

\begin{itemize}
  \item \textbf{Distributions are not yield.} The premium a fund collects and pays out comes with a cap on the fund's gains. In a strong market, the investor gives up upside in exchange for the income. The premium is not a dividend on top of the market return; part of it is the market return, converted into a different shape.
  \item \textbf{The fund is mostly an equity fund.} With a beta of around 0.5 to 0.7 and a down beta near 1, it will fall nearly as much as the market in a crash and rise much less in a rally. Compare it with a mix of an index fund and bills at the same beta, not with the index.
  \item \textbf{Look at the worst months, not the Sharpe ratio.} A smooth return history can hide a tail. Ask how the fund did in 2008 and in March 2020.
\end{itemize}

\textbf{Questions to ask a manager selling volatility:}

\begin{itemize}
  \item What is the strategy's beta, and its down beta? What would a beta-matched index position have returned?
  \item How far out of the money, and how short-dated, are the options sold? How is the strike reset?
  \item What is the largest loss the strategy could suffer in one day, and how was that figure derived?
  \item Is there leverage\index{leverage}, or margin that could be called in a spike?
\end{itemize}

\begin{callout}
\textbf{What Failed: selling volatility on margin, February 2018.} LJM Partners ran funds that used margin to sell out-of-the-money puts and calls on S\&P 500 futures. For years the strategy produced steady returns. In February 2018, when volatility spiked, the LJM-managed funds lost more than 80\%, over \$1 billion, in two trading days. The same week, the VelocityShares inverse VIX note described in Chapter~\ref{ch:10} lost 96\% in a day.

In 2021 the SEC charged LJM and two portfolio managers with misleading investors about the funds' risk management. According to the SEC, the firm had told investors it stress-tested for a worst-case daily loss of 20\% in its mutual fund, but that figure had been derived from a historical single-day loss of about 9\%, doubled and rounded up.

The episode illustrates this chapter's evidence in its sharpest form. The variance risk premium is real, and selling volatility earned it for years. But the premium is payment for exactly the kind of day that ended LJM. Leverage turned a strategy that should lose heavily in a spike into one that could not survive it.
\end{callout}

\section*{Exercises}
\addcontentsline{toc}{section}{Exercises}

\subsection*{Conceptual (all readers)}

\begin{enumerate}
  \item The index is at 4,000, the risk-free rate is 5\%, and one-month implied volatility is 18\%. Price a one-month at-the-money call and put by Black--Scholes and check put--call parity.
  \item Use put--call parity to show that a covered call and a cash-secured put with the same strike and expiry have the same payoff. Why do Cboe's BXM and PUT indices nonetheless have different returns?
  \item A covered call has a beta of 0.6 and has returned 7\% a year when the market returned 9\% and bills 2\%. What is its CAPM alpha? What return would a beta-matched portfolio of the index and bills have earned?
\end{enumerate}

\subsection*{Lab extensions (technical readers)}

\begin{enumerate}
  \setcounter{enumi}{3}
  \item Part A: in what share of months did the VIX exceed realized volatility? Compute the average gap separately for months with VIX above and below 25. Where is the premium larger? Then split the sample at 2008: how different is the premium before and after?
  \item Part B: rerun with \texttt{-{}-iv-shift -2} and \texttt{-4}. How much of the rebuilt put-write's alpha survives? Compare with Cboe's PUT if you downloaded it.
  \item Part C: compute the MPPM with risk aversion 1, 3 and 6. At what risk aversion does put-writing stop beating the market?
  \item Part D: what share of each strategy's excess return comes from beta? Compare the skewness of the put-write with that of the beta-matched market.
  \item Part D: the pure short variance swap has a high Sharpe ratio. Find its worst month. How much leverage could an investor have used before that month would have wiped out the account?
  \item A retiree is offered an option-income fund paying an 8\% distribution. Using this chapter's evidence, write a one-page explanation of what the fund is likely to deliver and what could go wrong.
\end{enumerate}
